\documentclass[11pt,a4paper]{article}
\usepackage[utf8]{inputenc}
\usepackage[T1]{fontenc}
\usepackage[italian]{babel}
\usepackage{mathptmx}
\usepackage[margin=2.4cm]{geometry}
\usepackage{enumitem}
\usepackage{longtable}
\usepackage{array}
\usepackage{hyperref}
\hypersetup{hidelinks}

\setlength{\parskip}{4pt}
\setlength{\parindent}{0pt}

\newcommand{\integrata}{\textbf{Integrata.}}
\newcommand{\parziale}{\textbf{Integrata in parte.}}
\newcommand{\nonintegrata}{\textbf{Non integrata.}}

\title{\textbf{Risposta alle revisioni}\\[4pt]
\large XXV CIM 2026 --- submission 159\\
\normalsize PythonGranularEngine: un ambiente dichiarativo per la granulazione
e la sua mappa sinottica}
\author{Giulio Romano De Mattia}
\date{31 agosto 2026}

\begin{document}
\maketitle

\section*{Premessa}

Ringrazio i due revisori per la lettura attenta e per le osservazioni
puntuali, che hanno orientato una revisione estesa dell'articolo. La
maggior parte delle indicazioni è stata recepita; le poche che non lo sono
state, o che lo sono solo in parte, sono elencate qui sotto con la
motivazione.

Tre osservazioni sono state condivise da entrambi i revisori e hanno guidato
il lavoro nel suo complesso:

\begin{enumerate}[leftmargin=*]
\item \textbf{Chiarezza e registro della prosa} (R1, punto 2; R2, nota
finale). Accolta integralmente. L'articolo è stato riscritto nelle parti
segnalate come oscure, e la revisione ha investito soprattutto abstract,
introduzione e conclusioni, che erano le sezioni più cariche di formulazioni
evocative. Tutte le locuzioni citate testualmente da R1 sono state eliminate
o riscritte in forma piana. Il criterio adottato: ogni frase deve dire un
fatto verificabile sul sistema o sulla figura che accompagna.

\item \textbf{Novità sopravvalutata / da dichiarare esplicitamente}
(R1, punto 1; R2, punto 2). Accolta con una concessione piena. La
rivendicazione che provocava l'obiezione è stata ritirata invece che
difesa: sono spariti dall'articolo tutti i claim di superiorità della
notazione testuale su altri ambienti, e tutte le asserzioni su \emph{come si
lavora} con Max, Pure Data, SuperCollider o Csound. Restano soltanto le
affermazioni positive su ciò che questo ambiente fa, dette senza termine di
paragone. Il motivo per cui non si risponde all'obiezione è che essa ha
ragione nella sua forma forte: la maggior parte dei compositori usa quegli
ambienti in modo di fatto dichiarativo, e nessuna affermazione su come altri
lavorano è difendibile in un articolo che non conduce uno studio d'uso.
Quanto resta di dichiarato come proprio è circoscritto e detto una volta
sola: il gate probabilistico che separa ampiezza e probabilità della
deviazione (\S\,\emph{La deviazione per grano}, con i due precursori più
vicini nominati esplicitamente, ICMS e EmissionControl2), e l'asse verticale
della \textsc{map} come posizione di lettura nel materiale (con il lineage
verso Truax 1988 Fig.~4, Roads e IRIN dichiarato in nota). Le tecniche
ereditate sono attribuite dove compaiono: tendency mask a Truax e CMask,
sincrono/quasi-sincrono/asincrono a Roads, \emph{granular sampling} a Lippe.

\item \textbf{Tempo reale e tempo differito} (R1, punto 1, nota a margine;
R2, punto 3). Accolta. R1 ha ragione: reale e differito sono proprietà del
\emph{motore}, non della descrizione, e il motore non era l'oggetto
dell'articolo. La discussione è stata rimossa: la sezione autonoma che
argomentava il tempo differito come premessa teorica non esiste più, e le
conclusioni non contengono più alcuna difesa del differito. Ciò che resta è
la sola affermazione difendibile e non comparativa, cioè che la \textsc{map}
è una rappresentazione dell'intera popolazione disponibile dopo il
rendering. Questo risponde insieme a R1 (che chiedeva di toglierla) e a R2
(che ne chiedeva una giustificazione migliore): la giustificazione non serve
più, perché il claim che la richiedeva è stato ritirato. Il taglio ha inoltre
liberato lo spazio usato per le aggiunte richieste altrove.
\end{enumerate}

\newpage
\section*{Revisore 1 --- osservazioni maggiori}

\begin{description}[leftmargin=0pt,style=unboxed]

\item[M1 --- Novità dell'elemento dichiarativo; equivalenza con i controlli
di una patch Max/PD; dire in quali contesti la notazione testuale convenga.]
\integrata\ Come descritto al punto 2 della premessa: la rivendicazione è
ritirata, non difesa. Non si argomenta più in quali contesti la notazione
testuale sia \emph{più efficace}, perché sarebbe di nuovo un'affermazione
comparativa non sostenuta da dati. Le ragioni della scelta testuale restano
nell'articolo dichiarate per quello che sono, cioè motivi di questo progetto
e non tesi generali: versionamento del processo con git, esposizione del
motore come API sotto un generatore di livello più alto, e separazione fra
generazione dei grani e rendering, che rende la lista di grani riusabile da
più moduli d'uscita.

\item[M2 --- La discussione tempo reale / tempo differito non serve; solo la
\textsc{map} è soggetta alla distinzione e il suo status è sopravvalutato.]
\integrata\ Per la prima metà, si veda il punto 3 della premessa. Per la
seconda: lo status della \textsc{map} è stato ridimensionato. Non è più
presentata come contributo teorico ma come strumento di lavoro, con la sua
funzione dichiarata in termini operativi (analisi, debugging, documentazione
del lavoro compositivo) e il suo lineage storico riconosciuto in nota. Le
conclusioni ne dichiarano inoltre un limite non colmato: l'altezza è affidata
al solo colore, che la stampa in bianco e nero e la visione cromatica atipica
non restituiscono.

\item[M3 --- Prosa a tratti incomprensibile; elenco di locuzioni; semplificare
a partire dall'introduzione.] \integrata\ Tutte le locuzioni citate dal
revisore sono state eliminate. Nessuna delle espressioni riportate fra
virgolette nella recensione (``la specifica si compila e si consuma'',
``l'artefatto normalizzato su cui operano le trasformazioni'', ``Nel
congelamento le due vie sono complementari'', ``come gli esempi successivi
dispiegano'', ``grani fedeli e grani devianti in proporzione componibile nel
tempo'', ``I primi tre angoli del quadro sono già a tema'', ``un controllo
che la tradizione tiene fuso o binario'', ``là subìto, qui scelto in quanto
abitabile'') compare più nel testo. L'introduzione è stata riscritta da capo
attorno a un problema enunciato in termini concreti: pochi secondi di
granulazione producono decine di migliaia di eventi, la event list che li
descrive è meno leggibile del file audio che ne deriva, e il file audio non
consente di risalire ai parametri che lo hanno generato.

\item[M4 --- Riportare i codici sotto le immagini corrispondenti.]
\integrata\ Ogni listato YAML è ora dentro il float della figura cui si
riferisce, sopra la \textsc{map} che ne è il rendering, sotto un'unica
didascalia. Listato e figura non possono più separarsi in impaginazione.

\item[M5 --- Ripensare e semplificare la terminologia.] \parziale\ Si veda
la nota dettagliata D19/D24 più sotto per il caso principale
(\texttt{dephase}), rinominato nel paper e nel software. Sul resto la
terminologia è stata riportata all'italiano corrente dove esisteva
(\emph{inviluppo} per \emph{envelope}) e i termini tecnici sono ora glossati
alla prima occorrenza.

\item[M6 --- Includere una specifica completa del linguaggio, con la
ramificazione di proprietà, sottoproprietà e valori.] \parziale\ La
richiesta è legittima e la specifica completa è stata effettivamente
prodotta, in forma di albero della grammatica generato automaticamente
dagli schemi del software. Non è però entrata nel PDF: sono 74 righe, cioè
poco meno di una colonna intera, e questa stessa revisione doveva
soddisfare la richiesta M7 dello stesso revisore, che chiedeva di
\emph{accorciare} l'articolo. Le due indicazioni non stavano insieme entro
il limite di 8 pagine, e ho dato la precedenza alla chiarezza e alla brevità
richieste da entrambi i revisori. Il paper rinvia perciò in nota alla
documentazione integrale della sintassi nel repository
(\texttt{docs/reference/yaml.md}), che è mantenuta col codice ed è la fonte
autorevole, mentre le chiavi necessarie a ciascun esempio sono introdotte
progressivamente nel testo.

\item[M7 --- L'articolo è troppo lungo per quello che propone.]
\integrata\ È stato tagliato in modo sostanziale, non cosmetico. Sono
sparite per intero la sezione autonoma sul rapporto con la tradizione
storica (compressa in poche righe nelle conclusioni) e la sezione sulle
implicazioni del tempo differito; il caso compositivo esteso è stato
eliminato. Lo spazio recuperato ha finanziato le aggiunte richieste dai
revisori (spettrogramma, tabella del jitter, didascalie estese) mantenendo
l'articolo entro le 8 pagine.

\item[M8 --- Equazioni e modelli probabilistici corretti ma inutili qui.]
\parziale\ Le equazioni sopravvissute sono tre e ciascuna sostiene
un'affermazione che il testo fa e che altrimenti resterebbe vaga:
l'interpolazione convessa fra griglia sincrona e asincrona, il modello della
tendency mask nella forma con centro e semiampiezza, e la sua estensione con
il fattore di gate. Le ultime due sono in particolare il modo più economico
di mostrare che l'estensione proposta si riduce esattamente al modello
classico quando la probabilità vale 1, che è il punto dell'articolo: senza
le due formule affiancate resterebbe un'affermazione da credere sulla
parola. Tutto il resto dell'apparato formale è stato rimosso.

\end{description}

\section*{Revisore 1 --- note dettagliate}

\begin{description}[leftmargin=0pt,style=unboxed]

\item[Abstract --- ``queryable IR'' non è chiaro.] \integrata\ L'abstract è
stato riscritto da capo e l'espressione non compare più. Dove serviva, la
rappresentazione intermedia è ora descritta per esteso e in termini di cosa
è: una lista di grani, ciascuno con il proprio insieme completo di
parametri, che alimenta insieme il rendering audio e la \textsc{map}.

\item[§1 --- ``distanza opaca'': in che senso dichiarativo, perché il
risultato non è leggibile.] \integrata\ La formulazione è stata sostituita
da un'enunciazione concreta: a quella scala ricondurre l'esito spettrale e
morfologico ai parametri che lo hanno prodotto richiede un controllo che la
sola specifica non fornisce, e senza un passaggio di analisi dei dati
verificare l'effetto di un parametro significa procedere per tentativi.

\item[§1 --- ``ciò che una specifica testuale non mostra'': una specifica
testuale può mostrare l'onset del grano.] \integrata\ Il revisore ha
ragione e la frase era sbagliata così com'era scritta. Il punto non è che il
testo non possa dirlo, ma che una event list di decine di migliaia di
eventi è, come documento, meno accessibile del file audio che ne deriva.
L'introduzione ora dice esattamente questo.

\item[§1 --- ``Tre tradizioni hanno attraversato questa distanza'':
paragrafo in cui il lettore si perde.] \integrata\ Il paragrafo è stato
eliminato. L'introduzione non contiene più una rassegna di tradizioni.

\item[§1 --- SuperCollider ``dichiara l'esito e non la procedura'': dubbio
fondato.] \integrata\ L'affermazione è stata rimossa, insieme a tutte le
altre asserzioni su come si lavora con altri ambienti (si veda M1).

\item[§2 --- Figura 1: la legenda ``pitch (cents)'' non è usata.]
\integrata\ La didascalia della prima figura ora dichiara esplicitamente
che il colore codifica la trasposizione, ``qui nulla'', e la colorbar è
presentata nel cappello di sezione insieme agli altri elementi di lettura
della \textsc{map}, prima che la figura compaia.

\item[§2 --- Wrapping fine/inizio: buffer circolare? click? campioni a
cavallo?] \integrata\ L'assunzione è ora esplicita nel testo: l'indice di
lettura è preso modulo la durata del file, il buffer si comporta come una
tabella letta ciclicamente, e il grano a cavallo della giunzione la
attraversa senza interrompersi. Il salto verticale visibile nella figura è
una discontinuità della posizione di lettura, non un reset del fasore, che
continua ad accumulare fase. Sulla questione del click: la finestra
applicata a ciascun grano porta a zero gli estremi, quindi la giunzione non
produce discontinuità di ampiezza.

\item[§2 --- Il verso di lettura cambia: perché, e non dovrebbe essere un
parametro a parte?] \integrata\ Lo è, e ora il testo lo dice. Dove la
\texttt{speed\_ratio} è negativa si invertono insieme il verso di
percorrenza del buffer e quello di lettura interno ai grani; questo
accoppiamento è il comportamento di default e vale finché non si dichiara
\texttt{grain.reverse}, chiave che rende i due versi dichiarabili
separatamente.

\item[§2 --- ``non visualizzabile con forma d'onda o sonogramma'': sembra
cinematica di base.] \integrata\ La frase è stata ridimensionata. Non si
afferma più che il comportamento sia incomprensibile, ma il fatto
verificabile: la posizione di lettura è un asse della \textsc{map}, e non è
un asse della forma d'onda né dello spettrogramma.

\item[§2 --- Nota 7 sembra un esponente; \texttt{fill\_factor} è l'overlap e
va nel testo.] \integrata\ La spaziatura del richiamo di nota è stata
corretta. La sovrapposizione è ora nominata nel corpo del testo e nella
didascalia della prima figura, dove si legge la ``sovrapposizione a
fattore~2''; la definizione formale di \emph{fill factor} di Roads resta in
nota come riferimento.

\item[§2 --- Figura 2 illeggibile per densità di grani troppo alta; va
spostata dopo.] \parziale\ La leggibilità è stata risolta spostando le due
lenti di ingrandimento su punti a bassa densità della stessa figura: dentro
la lente i singoli onset sono ora distinguibili, ed è lì che si legge il
confronto fra griglia regolare e griglia irregolare. La densità dell'esempio
non è stata abbassata perché l'esempio è stato scelto per il suo risultato
sonoro, che fa parte dell'archivio audio pubblicato; abbassarla avrebbe
prodotto una figura più chiara ma un esempio diverso da quello che si
ascolta. La posizione della figura è stata rivista in coerenza con la nuova
sequenza degli esempi.

\item[§2 --- Perché IOT ha una linea superiore? Meglio IOT\_average.]
\integrata\ La notazione è ora $\mathrm{IOT}_{\mathrm{avg}}$ in tutte le
occorrenze, testo, equazione e note.

\item[§2 --- ``in ascissa il tempo del brano, in ordinata la posizione di
lettura\ldots'' va nella didascalia.] \integrata\ Le descrizioni degli assi
e della lettura delle figure sono state spostate nelle rispettive
didascalie, in tutte le figure dell'articolo. Il corpo del testo rinvia alla
figura invece di descriverla.

\item[§2 --- ``un treno periodico di grani ha uno spettro a righe'':
servirebbero illustrazioni spettrali, non temporali.] \integrata\ È stato
aggiunto uno spettrogramma sotto la \textsc{map} della figura sulla griglia
temporale, e la didascalia lega i due pannelli: dove la griglia è ancora
sincrona e la densità alta si vedono le righe spaziate a multipli della
densità; dove \texttt{distribution} sale verso 1 le righe si smerigliano.
L'affermazione non è più solo asserita.

\item[§2 --- ``Nel congelamento le due vie sono complementari'': congelamento
di cosa?] \integrata\ La frase è stata eliminata.

\item[§2 --- Il listato 3 non è referenziato né commentato.] \integrata\
Ogni listato è ora dentro il float della propria figura e commentato dalla
didascalia corrispondente; non esistono più listati privi di riferimento nel
testo.

\item[§2 --- ``envelope'' è \emph{inviluppo} in italiano.] \integrata\
Sostituzione sistematica in tutto l'articolo. Il termine inglese resta solo
dove nomina la chiave YAML \texttt{Envelope}, cioè un identificatore del
linguaggio e non una parola del testo.

\item[§2 --- Manca il numero della figura in ``etichettati (a) e (b) nella
figura''.] \integrata\ Il riferimento è ora esplicito.

\item[§2 --- \texttt{dephase.pointer} è un nome fuorviante; forse
\texttt{offset\_probability} sotto \texttt{pointer}; l'implementazione è
barocca.] \parziale\ La sostanza dell'osservazione è accolta e ha prodotto
una rinomina non solo nell'articolo ma nel software, rilasciata come
versione di rottura senza alias di retrocompatibilità: la chiave si chiama
ora \texttt{deviation\_probability}, che dice quello che il parametro è, cioè
una probabilità. Non è stata accolta la struttura proposta, cioè annidarla
sotto \texttt{pointer}: la stessa chiave ha anche una forma globale in cui
si applica a tutti i parametri insieme, e nidificarla sotto uno di essi
cancellerebbe quella forma. Non è stato adottato nemmeno il nome
\texttt{jitter} suggerito altrove, perché nel software \emph{jitter} nomina
già l'ampiezza della deviazione: userebbe la stessa parola per i due fattori
che l'articolo distingue.

\item[§2 --- ``le due morfologie non potrebbero essere più diverse'':
colloquiale; spiegazione floreale.] \integrata\ Il passo è stato riscritto
per fatti: la modulazione dell'ampiezza produce un allargamento graduale
della nuvola, la modulazione della probabilità agisce come interpolazione
fra due popolazioni, quelli che seguono la traiettoria e quelli che
deviano.

\item[§2 --- Eq. 2: cos'è $v_n$?] \integrata\ Tutti i simboli sono ora
definiti: $v_n$ è il valore assegnato al grano $n$ al tempo normalizzato
$\tau_n$, $c$ la traiettoria centrale, $\rho$ l'ampiezza della deviazione,
$\xi_n$ il campionamento casuale, $g_n$ il gate.

\item[§2 --- ``Il caso (a) è questo modello all'opera'', con $c(\tau_n)$
costante.] \integrata\ Precisato: il pannello (a) mostra il modello nella
sua forma classica, con traiettoria costante e deviazione crescente nel
tempo.

\item[§2 --- ``Ampiezza e probabilità sono due inviluppi ortogonali, (a)
muove la prima a gate aperto\ldots'': non si capisce.] \integrata\ Riscritto
dicendo esplicitamente, per ciascuno dei due stream, quale inviluppo si
muove e quale resta fermo.

\item[§2 --- ``dephase'' non è deviazione, forse \emph{jitter}.] \parziale\
Si veda sopra: la sostanza è accolta con la rinomina, il nome proposto no,
per la ragione detta.

\item[§2 --- Figura 4: l'andamento ``cubico'' 20--80\% non si vede; e
deviazione = 100 di cosa?] \integrata\ Il revisore aveva ragione su
entrambi i punti, ed erano due errori. La didascalia diceva ``cubico''
mentre l'esempio dichiara un inviluppo a gradini: corretta in ``quattro
gradini''. Quanto all'unità, la curva etichettata con il simbolo di
percentuale non misura l'ampiezza della deviazione ma la probabilità che essa
venga applicata: la didascalia ora lo dichiara esplicitamente e distingue le
due curve, rinviando alla figura precedente per l'ampiezza.

\item[§2 --- Perché Mid-Side e Blumlein a questo livello di astrazione?
Perché non estendere oltre la circonferenza unitaria?] \parziale\ La
citazione è stata mantenuta ma spostata in nota e ridotta al ruolo di
riferimento della codifica, non di argomento. La nota ora dice esattamente
ciò che il revisore osserva: $\theta$ è un azimuth astratto, Left e Right
sono la stessa coppia ruotata di 45 gradi e non l'unico esito possibile, e
ruotando con un'altra matrice si ottengono microfoni virtuali verso $n$
altoparlanti qualsiasi. Il limite alla circonferenza è dichiarato come
limite del motore, che resta agnostico sul numero di canali d'uscita: lo
stereo è il default, non un vincolo dell'architettura. L'estensione a uno
spazio di dimensione maggiore non è stata implementata e non viene
promessa.

\item[§2 --- \texttt{pan\_range}: jitter sul pan o apertura del fronte? Il
nome è infelice, e ``come ogni parametro'' non torna.] \integrata\ La regola
generale è ora enunciata come regola prima degli esempi, nel cappello della
sezione, invece di comparire tardi e già applicata: quasi ogni parametro si
declina in valore base, banda di deviazione e probabilità che la deviazione
si applichi. Il caso del pan è ora descritto esplicitamente nei suoi due
livelli annidati: la strategia del blocco \texttt{voices} assegna a ciascuna
voce un proprio angolo, distribuendo le voci nel campo, mentre
\texttt{pan\_range} aggiunge a ogni grano della singola voce una deviazione
attorno a quell'angolo. È dunque il secondo dei due sensi che il revisore
distingueva.

\item[§3 --- ``La svolta è documentabile negli atti del 1993'': il real-time
granulare è precedente (Riverrun, 1986); i riferimenti sono troppo centrati
su CIM.] \integrata\ La sezione che conteneva quella periodizzazione è stata
rimossa nella revisione, quindi l'affermazione contestata non esiste più.
Sul secondo punto: la bibliografia è stata ampliata fuori dagli atti CIM e
comprende ora Roads (\emph{Microsound} e i lavori sui pulsar), Truax 1988,
Bartetzki su CMask, Dutilleux e altri sul trattamento canonico della
granulazione, Vaggione, Blumlein e Gerzon per la parte spaziale, Seeger per
la distinzione fra notazione prescrittiva e descrittiva.

\item[§3 --- ``dove la specifica resta interrogabile'': servono casi
concreti; cosa fa lo YAML che gli oggetti Max non fanno.] \integrata\ Non si
risponde al confronto perché il confronto è stato ritirato (si veda M1). La
formulazione contestata è stata eliminata.

\item[§3 --- ``esporre come asse dichiarativo continuo\ldots Due casi la
circoscrivono'': referente ignoto, paragrafo fumoso.] \integrata\ Il
paragrafo è stato eliminato con la sezione che lo conteneva. Quanto resta
del contenuto è detto in forma piana dove serve, cioè nella sezione sulla
deviazione, con i due precursori nominati per nome invece che allusi.

\item[§4 --- ``perché l'architettura tiene il ritorno dell'ascolto
abbastanza vicino alla scrittura da fare del ciclo uno spazio di lavoro'':
non chiaro.] \integrata\ La frase è stata eliminata. Le conclusioni sono
state riscritte e dicono ora tre cose verificabili: cosa fa il gate
probabilistico, che la specifica resta un documento versionabile e
rigenerabile dato il seed dichiarato, e che la \textsc{map} mostra il
risultato prima del riascolto. Seguono tre limiti dichiarati e non colmati.

\end{description}

\section*{Revisore 2}

\begin{description}[leftmargin=0pt,style=unboxed]

\item[R2.1 --- L'utilità della \textsc{map} è asserita più che valutata;
negli esempi complessi l'informazione diventa densa e potenzialmente
difficile da interpretare.] \integrata\ La \textsc{map} è ora presentata
prima di comparire, nel cappello di sezione, con la spiegazione esplicita di
come si legge: assi, forma d'onda laterale, glifo del grano, orientamento
della testa, profilo della finestra, colore per la trasposizione, pannello
degli inviluppi, etichetta di stream. La motivazione è dichiarata in termini
d'uso: è negli esempi fitti, con più stream e più parametri attivi, che la
vista d'insieme serve di più, per il test e il debugging e per capire dove
intervenire puntualmente su una popolazione di grani. Nell'articolo si è
scelto di mostrare un solo stream per figura proprio per non stressare la
leggibilità della figura stampata. Il limite di leggibilità che il revisore
intuisce è inoltre dichiarato nelle conclusioni per il caso più critico,
quello dell'altezza affidata al solo colore. Una valutazione condotta con
altri compositori è riconosciuta come lavoro non svolto e rimandato, e non
viene spacciata per fatta.

\item[R2.2 --- La distinzione fra tecniche ereditate e aspetti originali
resta dispersa; servirebbe un riepilogo esplicito dei contributi.]
\parziale\ La distinzione è stata resa esplicita, ma senza aggiungere un
riquadro dedicato. Il motivo è di spazio: R1 chiedeva contemporaneamente di
accorciare l'articolo, e ho preferito investire lo spazio disponibile nella
chiarezza del corpo. Il risultato è che ciò che è ereditato è ora attribuito
per nome nel punto in cui compare (tendency mask a Truax e a CMask, i tre
regimi di sincronia a Roads, il \emph{granular sampling} a Lippe, le
maschere di tendenza visualizzate a Truax 1988 Fig.~4), e ciò che è proprio è
detto una volta sola e circoscritto, nel paragrafo che segue le due
equazioni: il gate probabilistico come asse dichiarabile e componibile nel
tempo, con i due precursori più vicini nominati e la differenza rispetto a
ciascuno spiegata. Le conclusioni ricapitolano gli stessi due punti in tre
paragrafi. La rivendicazione complessiva è comunque più piccola di quella
della versione sottomessa, come richiesto da R1.

\item[R2.3 --- La giustificazione del rendering in tempo differito convince
solo in parte: perché la stessa rappresentazione dichiarativa e la stessa
analisi visiva non potrebbero stare in un sistema real-time o ibrido?]
\integrata\ Si veda il punto 3 della premessa. La risposta è che non c'è
ragione per cui non potrebbero, e infatti l'articolo non lo afferma più: la
pretesa che il differito fosse condizione necessaria dell'approccio è stata
ritirata insieme alla sezione che la sosteneva. Questo converge con
l'osservazione di R1, che chiedeva di togliere la stessa discussione per una
ragione affine.

\item[R2.4 --- Registro incoerente, alternanza fra colloquiale ed elaborato,
virtuosismo linguistico sopra la precisione.] \integrata\ Si veda M3: è la
stessa riscrittura, condotta sull'intero articolo con priorità ad abstract,
introduzione e conclusioni.

\end{description}

\section*{Nota finale}

Oltre a quanto richiesto dai revisori, la versione camera-ready contiene due
cambiamenti che vale la pena segnalare al Comitato. Il primo è il nome del
sistema: l'acronimo usato nella submission è stato abbandonato, e il sistema
è chiamato ovunque con il nome del suo repository pubblico,
PythonGranularEngine. Il secondo è la riproducibilità degli esempi: il
software è stato esteso con un seeding deterministico dichiarabile nella
specifica, e gli esempi dell'articolo dichiarano il proprio seed, sicché le
figure stampate e l'audio dell'archivio sono ciò che si riottiene eseguendo
i listati pubblicati.

\end{document}
